%%
%% This is file `example_es.tex',
%% generated with the docstrip utility.
%%
%% The original source files were:
%%
%% coppe.dtx  (with options: `examplees')
%% 
%% This is one of the demonstration documents of the `coppe' document class:
%% the five main-language demos, the PDF/A build of the full example and the
%% side-by-side cover sheet.  It is generated from coppe.dtx by docstrip; edit
%% the .dtx source, not this file.
%% 
%% Copyright (C) 2008-2026 Vicente Helano F. Batista, George O. Ainsworth Jr.,
%% Geraldo Xexéo and Eduardo Mangeli. Released under the GNU General Public
%% License version 3 (see the COPYING file).
%% 
%% \CheckSum{0}
%% \CharacterTable
%%  {Upper-case    \A\B\C\D\E\F\G\H\I\J\K\L\M\N\O\P\Q\R\S\T\U\V\W\X\Y\Z
%%   Lower-case    \a\b\c\d\e\f\g\h\i\j\k\l\m\n\o\p\q\r\s\t\u\v\w\x\y\z
%%   Digits        \0\1\2\3\4\5\6\7\8\9
%%   Exclamation   \!     Double quote  \"     Hash (number) \#
%%   Dollar        \$     Percent       \%     Ampersand     \&
%%   Acute accent  \'     Left paren    \(     Right paren   \)
%%   Asterisk      \*     Plus          \+     Comma         \,
%%   Minus         \-     Point         \.     Solidus       \/
%%   Colon         \:     Semicolon     \;     Less than     \<
%%   Equals        \=     Greater than  \>     Question mark \?
%%   Commercial at \@     Left bracket  \[     Backslash     \\
%%   Right bracket \]     Circumflex    \^     Underscore    \_
%%   Grave accent  \`     Left brace    \{     Vertical bar  \|
%%   Right brace   \}     Tilde         \~}
%%

\documentclass[spanish,dsc]{coppe}
\addbibresource{example.bib}
\begin{document}
  \title{Demonstração da classe CoppeTeX 4.1 em português}
  \foreigntitle{CoppeTeX 4.1 class demonstration in English}
  \titlein{spanish}{Demostración de la clase CoppeTeX 4.1 en espa\~nol}
  \author{Nombre del}{Autor}
  \advisor{Primer}{Director}{D.Sc.}{UFRJ}
  \examiner{Primer Examinador}{D.Sc.}{UFRJ}
  \examiner{Segundo Examinador}{D.Sc.}{UFRJ}
  \department{PESC}
  \date{05}{2026}
  \keyword{Multilenguaje}
  \keyword{Demostraci\'on}
  %% palavras-chave do resumo em idioma estrangeiro (3.1.2.1.5, Anexo F)
  \foreignkeyword{Multilingual}
  \foreignkeyword{Demonstration}
  %% palavras-chave do terceiro resumo, em português (Norma COPPE, secao 3)
  \braziliankeyword{Multilíngue}
  \braziliankeyword{Demonstração}
  \maketitle
  \frontmatter
  \begin{abstract}
    Este es el resumen principal en espa\~nol. La tesis est\'a escrita
    en espa\~nol como idioma principal, con la portada conservada en
    portugu\'es por convenci\'on institucional de la UFRJ. Lorem ipsum
    dolor sit amet, consectetur adipiscing elit.
  \end{abstract}
  \begin{foreignabstract}
    This is the English foreign abstract. Mandatory in CoppeTeX whichever
    the main language is, for international readability.
  \end{foreignabstract}
  \begin{brazilianabstract}
    Este \'e o resumo em portugu\^es, oferecido pelo ambiente
    \texttt{brazilianabstract} para a leitura da banca brasileira quando o
    idioma principal n\~ao \'e o portugu\^es.
  \end{brazilianabstract}
  % Listas pretextuales: solo el sumario es obligatorio (Manual UFRJ
  % 3.1.2.1.6). Las listas de figuras, tablas, abreviaturas y símbolos son
  % opcionales (3.1.2.2); las de cuadros, programas y algoritmos aparecen
  % cuando el trabajo tiene uno (Norma COPPE, §10). Esta demostración no
  % tiene ilustraciones, por eso solo imprime el sumario.
  \tableofcontents
  \mainmatter
  \chapter{Introducci\'on}
  Este cap\'itulo demuestra el cuerpo en espa\~nol, mientras la portada
  y la folha-de-rosto permanecen en portugu\'es (plantilla institucional
  UFRJ).
  \section{El archivo de depósito: PDF/A}
  El depósito en la UFRJ es exclusivamente digital desde la Resolución CEPG
  n.~246/2023, y el Manual de la UFRJ (2.2d) exige el archivo final en
  \textbf{PDF/A}, el formato de archivo a largo plazo. PDF/A es el PDF
  restringido por la norma ISO~19005 para la preservación: todas las fuentes van
  incrustadas, los colores declaran un perfil (sRGB) que los hace independientes
  del equipo, los metadatos van en un bloque XMP que los repositorios leen, y
  nada depende del exterior --- sin contenido externo, JavaScript ni cifrado.
  Así el archivo se abre igual dentro de décadas, en cualquier computadora.

  La clase usa la variante \textbf{PDF/A-2b}: la parte 2 admite transparencia,
  que la parte 1 prohíbe y que producen las cajas de color y los gráficos, y el
  nivel \emph{b} garantiza la reproducción visual fiel. Basta la opción
  \texttt{pdfa} en \verb|\documentclass[spanish,dsc,pdfa]{coppe}|: la clase carga
  \texttt{pdfx}, arma los metadatos a partir del título, el autor y las palabras
  clave, e incrusta el perfil de color. Actívela temprano, como hace
  \texttt{example.tex}: un PDF o una imagen fuera de la norma aparecen el día en
  que entran, y no la víspera del depósito. Para verificar, use el validador
  veraPDF, perfil PDF/A-2b.

  \section{Flujos de escritura}
  pdf\LaTeX{} mantiene como máximo dieciséis archivos abiertos para escritura,
  y cada lista, índice o glosario ocupa uno. Por eso la clase carga
  \texttt{morewrites}. La opción \texttt{semmorewrites} lo desactiva y ahorra
  cerca de medio segundo por pasada, pero solo sirve para un trabajo sin índice
  alfabético, sin varios índices y sin glosario.
\end{document}
%% 
%%
%% End of file `example_es.tex'.
